% GENERATED FILE. Edit canonical lessons, then run npm run generate:books.
% canonical-source-hash: fnv1a64:7138bbfd478dd774
% canonical-lessons: ES-C611-condicion, ES-C611-siempre-que, ES-C611-a-menos-que, ES-C611-con-tal-de-que, ES-C611-en-caso-de-que

\chapter{Con una condición}
\label{ch:es-con-una-condicion}
\hlchaptermodality{611}

\begin{chapteropening}
\textbf{By the end of this chapter:} \emph{I can set a condition with as long as, unless and provided that.}

Complete the last lesson of chapter 611: I can set a condition with as long as, unless and provided that.
\end{chapteropening}

\section[la condición]{la condición — the condition}

\label{lesson:ES-C611-condicion}



\begin{quote}
Before the new one: say the Spanish for the immigration, then the Spanish for the living together.
\end{quote}

\subsection*{You'll want to know: la condición}
\textbf{la condición} — \textquotedblleft{}the condition\textquotedblright{}.

\begin{quote}
Te ayudo con una condición.
\end{quote}

\begin{quote}
\emph{I will help you on one condition.}
\end{quote}

\emph{Las condiciones de trabajo} — working conditions.

\begin{cousinweb}
\emph{condición} is a learned borrowing from Latin \emph{condicio} 'agreement, terms', built on \emph{condicere} 'to agree together', from \emph{dicere} 'to say' (EtymDB).
\end{cousinweb}

\begin{grammarlens}[title={What you've built}]
One more word for telling, explaining and describing at B1.
\end{grammarlens}

\subsection*{Your turn}
\begin{itemize}
\raggedright
  \item \emph{Say it:} \emph{la condición}
  \item \emph{Say it:} \emph{la condición}, once more
  \item \emph{Say it:} a sentence of your own with \emph{la condición}
  \item \emph{Recall:} say the Spanish for the immigration, then the Spanish for the living together, then say \emph{la condición} again
\end{itemize}

\begin{tcolorbox}[breakable,colback=teal!4,colframe=teal!35!black,title={Before you move on}]
What does la condición mean? (\textquotedblleft{}The condition\textquotedblright{}.) Say it once more.
\end{tcolorbox}
\section[siempre que]{siempre que — as long as; whenever}

\label{lesson:ES-C611-siempre-que}



\begin{quote}
Before the new one: say the Spanish for the living together, then the Spanish for the condition.
\end{quote}

\subsection*{You'll want to know: siempre que}
\textbf{siempre que} — \textquotedblleft{}as long as; whenever\textquotedblright{}.

\begin{quote}
Puedes ir siempre que vuelvas pronto.
\end{quote}

\begin{quote}
\emph{You can go as long as you come back early.}
\end{quote}

With the subjunctive it sets a condition; with the indicative it means whenever.

\begin{cousinweb}
\emph{siempre que} is built on \emph{siempre}, inherited from Latin \emph{semper} 'always', and \emph{que} (EtymDB).
\end{cousinweb}

\begin{grammarlens}[title={What you've built}]
One more word for telling, explaining and describing at B1.
\end{grammarlens}

\subsection*{Your turn}
\begin{itemize}
\raggedright
  \item \emph{Say it:} \emph{siempre que}
  \item \emph{Say it:} \emph{siempre que}, once more
  \item \emph{Say it:} a sentence of your own with \emph{siempre que}
  \item \emph{Recall:} say the Spanish for the living together, then the Spanish for the condition, then say \emph{siempre que} again
\end{itemize}

\begin{tcolorbox}[breakable,colback=teal!4,colframe=teal!35!black,title={Before you move on}]
What does siempre que mean? (\textquotedblleft{}As long as; whenever\textquotedblright{}.) Say it once more.
\end{tcolorbox}
\section[a menos que]{a menos que — unless}

\label{lesson:ES-C611-a-menos-que}



\begin{quote}
Before the new one: say the Spanish for the condition, then the Spanish for as long as.
\end{quote}

\subsection*{You'll want to know: a menos que}
\textbf{a menos que} — \textquotedblleft{}unless\textquotedblright{}.

\begin{quote}
Voy a la playa, a menos que llueva.
\end{quote}

\begin{quote}
\emph{I am going to the beach, unless it rains.}
\end{quote}

Always followed by the subjunctive.

\begin{cousinweb}
\emph{a menos que} is built on \emph{menos}, inherited from Latin \emph{minus} 'less' (EtymDB).
\end{cousinweb}

\begin{grammarlens}[title={What you've built}]
One more word for telling, explaining and describing at B1.
\end{grammarlens}

\subsection*{Your turn}
\begin{itemize}
\raggedright
  \item \emph{Say it:} \emph{a menos que}
  \item \emph{Say it:} \emph{a menos que}, once more
  \item \emph{Say it:} a sentence of your own with \emph{a menos que}
  \item \emph{Recall:} say the Spanish for the condition, then the Spanish for as long as, then say \emph{a menos que} again
\end{itemize}

\begin{tcolorbox}[breakable,colback=teal!4,colframe=teal!35!black,title={Before you move on}]
What does a menos que mean? (\textquotedblleft{}Unless\textquotedblright{}.) Say it once more.
\end{tcolorbox}
\section[con tal de que]{con tal de que — provided that}

\label{lesson:ES-C611-con-tal-de-que}



\begin{quote}
Before the new one: say the Spanish for as long as, then the Spanish for unless.
\end{quote}

\subsection*{You'll want to know: con tal de que}
\textbf{con tal de que} — \textquotedblleft{}provided that\textquotedblright{}.

\begin{quote}
Te lo presto con tal de que me lo devuelvas.
\end{quote}

\begin{quote}
\emph{I will lend it to you provided that you give it back.}
\end{quote}

Followed by the subjunctive; \emph{con tal de} + infinitive also works.

\begin{cousinweb}
\emph{con tal de que} is built on \emph{tal}, inherited from Latin \emph{talis} 'such, of that kind' (EtymDB).
\end{cousinweb}

\begin{grammarlens}[title={What you've built}]
One more word for telling, explaining and describing at B1.
\end{grammarlens}

\subsection*{Your turn}
\begin{itemize}
\raggedright
  \item \emph{Say it:} \emph{con tal de que}
  \item \emph{Say it:} \emph{con tal de que}, once more
  \item \emph{Say it:} a sentence of your own with \emph{con tal de que}
  \item \emph{Recall:} say the Spanish for as long as, then the Spanish for unless, then say \emph{con tal de que} again
\end{itemize}

\begin{tcolorbox}[breakable,colback=teal!4,colframe=teal!35!black,title={Before you move on}]
What does con tal de que mean? (\textquotedblleft{}Provided that\textquotedblright{}.) Say it once more.
\end{tcolorbox}
\section[en caso de que]{en caso de que — in case; if}

\label{lesson:ES-C611-en-caso-de-que}



\begin{quote}
Before the new one: say the Spanish for unless, then the Spanish for provided that.
\end{quote}

\subsection*{You'll want to know: en caso de que}
\textbf{en caso de que} — \textquotedblleft{}in case; if\textquotedblright{}.

\begin{quote}
En caso de que llegues tarde, llámame.
\end{quote}

\begin{quote}
\emph{If you are late, call me.}
\end{quote}

Followed by the subjunctive; \emph{en caso de} + noun: \emph{en caso de incendio} — in case of fire.

\begin{cousinweb}
\emph{en caso de que} is built on \emph{caso}, a learned borrowing from Latin \emph{casus} 'a falling, an event', from \emph{cadere} 'to fall' (EtymDB).
\end{cousinweb}

\begin{grammarlens}[title={What you've built}]
One more word for telling, explaining and describing at B1.
\end{grammarlens}

\subsection*{Your turn}
\begin{itemize}
\raggedright
  \item \emph{Say it:} \emph{en caso de que}
  \item \emph{Say it:} \emph{en caso de que}, once more
  \item \emph{Say it:} a sentence of your own with \emph{en caso de que}
  \item \emph{Recall:} say the Spanish for unless, then the Spanish for provided that, then say \emph{en caso de que} again
\end{itemize}

\begin{tcolorbox}[breakable,colback=teal!4,colframe=teal!35!black,title={Before you move on}]
What does en caso de que mean? (\textquotedblleft{}In case; if\textquotedblright{}.) Say it once more.
\end{tcolorbox}
